\chapter{LITERATURE SURVEY}
\label{ch:literature-survey}

\section{Overview}
Zero-day intrusion detection focuses on identifying previously unseen attacks
without relying on signature-based databases, which are limited to known threat
patterns. In modern cybersecurity environments, where attack techniques evolve
rapidly, traditional methods struggle due to the absence of labeled data for new
threats. As a result, recent research has shifted towards learning normal system
behavior using self-supervised and unsupervised approaches, where anomalies are
detected as deviations from learned patterns rather than predefined signatures.

Graph-based representation learning has emerged as an effective approach for
modeling structured interactions between system entities such as processes,
files, and network connections. These methods capture complex relationships and
enable anomaly detection through reconstruction error. Additionally, machine
learning techniques for zero-day detection emphasize generalizing normal behavior
to identify unseen attacks. To ensure long-term effectiveness, continual learning
and dynamic model updates allow systems to adapt to evolving environments, while
alert correlation improves interpretability by converting anomaly scores into
meaningful and actionable insights.

\section{Self-Supervised Graph Representation Learning}
Guerra et al. propose a self-supervised framework for learning graph
representations in intrusion detection systems [2]. Their approach models network
traffic as graphs, where nodes represent entities and edges capture interactions
between them. The model employs Graph Neural Networks (GNNs) combined with
transformer-based encoders to learn both local and global structural patterns.

The training process is based on reconstructing graph structures, allowing the
model to learn normal behavior without requiring labeled attack data. During
inference, anomalies are detected through high reconstruction error, which
indicates deviations from learned patterns. The study demonstrates that graph
structure alone can provide strong anomaly signals, achieving high performance
even in the absence of supervised labels.

This work directly influences our system design, where behavioral data is
represented as graphs and processed using autoencoder-based models. The concept
of using reconstruction error as an anomaly score is adopted as a core detection
mechanism in our pipeline.

\section{Time-Aware Heterogeneous Hypergraph Learning}
Hayat et al. introduce a self-supervised framework for learning from
time-aware heterogeneous hypergraphs [4]. Unlike traditional graphs, hypergraphs
can model complex relationships involving multiple entities simultaneously,
making them suitable for representing interactions among processes, files, and
network connections.

Their model incorporates temporal information by analyzing sequences of events
over time windows, enabling the capture of evolving behavioral patterns. The
framework learns embeddings that preserve both structural and temporal features,
resulting in improved performance for dynamic graph-level classification tasks.

This work supports two key design decisions in our system. First, it justifies
the use of heterogeneous entities such as processes, files, and sockets within a
unified graph structure. Second, it emphasizes the importance of time-window-based
graph construction, which we adopt to capture temporal evolution in system
behavior.

\section{Zero-Day Detection Using Machine Learning}
Dai et al. focus on detecting zero-day attacks using machine learning techniques
that generalize to unseen data [1]. Their model is designed to identify anomalies
without relying on predefined attack signatures, making it effective in detecting
previously unknown threats.

The study highlights the importance of learning generalized representations of
normal behavior rather than memorizing known attack patterns. By training on
representative datasets, the model can distinguish between benign and anomalous
activities in unseen environments.

This aligns closely with our system’s self-supervised approach, where training is
performed exclusively on normal data. Anomalies are detected as deviations from
learned patterns, enabling effective zero-day detection without requiring labeled
attack datasets.

\section{Alert Correlation and Detection Pipelines}
Lanigan et al. address the challenge of alert correlation in intrusion detection
systems [5]. Traditional systems often generate large volumes of isolated alerts,
leading to difficulty in identifying meaningful threats. Their work proposes
methods to correlate alerts based on context, severity, and relationships.

By structuring alerts and enabling correlation, the system can identify attack
patterns and reduce false positives. This improves both interpretability and
response efficiency in real-world security operations.

This study informs the design of our alerting module, where anomaly scores are
converted into structured alerts with severity levels and contextual metadata.
This enables downstream analysis and integration with security monitoring
systems.

\section{Continual Learning for Intrusion Detection}
Guo et al. propose continual learning techniques to address the challenge of
evolving network behavior [3]. In real-world environments, normal activity
patterns change over time, causing static models to become outdated.

Their framework allows models to incrementally learn from new data while
retaining previously acquired knowledge. This is achieved through strategies that
mitigate catastrophic forgetting, ensuring that the model remains effective over
time.

This concept is highly relevant to our system, as it highlights the need for
adaptability. Our modular design allows for future integration of continual
learning techniques, enabling the system to update itself with new behavioral
patterns without losing prior knowledge.

\section{Dynamic Model Updates Using Threat Intelligence}
Lin et al. emphasize the role of cyber threat intelligence in enabling dynamic
updates for machine learning-based intrusion detection systems [6]. Their work
proposes integrating external threat data to continuously refine and improve
model performance.

The study highlights the importance of balancing adaptability with stability,
ensuring that model updates enhance detection capabilities without introducing
instability. By incorporating new threat intelligence, the system remains
relevant in rapidly evolving environments.

This work supports the design of our system’s update mechanism, where models can
be retrained or fine-tuned using new data. It reinforces the importance of
keeping detection systems aligned with current threat landscapes.

\section{Analysis of Continual Learning Models}
Prasath et al. provide a comparative analysis of various continual learning
approaches applied to intrusion detection systems [7]. Their study evaluates
different methods based on performance, scalability, and resistance to
catastrophic forgetting.

The authors highlight that no single approach is universally optimal, and the
choice of method depends on system requirements such as data availability and
computational constraints. They also emphasize the trade-off between accuracy
and efficiency in real-time detection systems.

This analysis provides valuable insights for future enhancements of our system.
It enables informed selection of continual learning strategies, ensuring that
the system remains both efficient and adaptable in dynamic environments.

\section{SUMMARY}
The selected literature establishes a comprehensive foundation for the proposed
zero-day intrusion detection system. Self-supervised graph learning enables
effective anomaly detection without requiring labeled attack data, which is
critical for identifying zero-day threats. Graph-based representations capture
structural relationships between system entities, while temporal modeling ensures
that evolving behavior is effectively analyzed.

Machine learning approaches for zero-day detection reinforce the importance of
learning normal behavior rather than relying on known attack signatures. Alert
correlation techniques improve interpretability and usability by transforming
raw anomaly scores into meaningful insights. Furthermore, continual learning and
dynamic model update strategies ensure that the system remains effective over
time, adapting to changing environments and emerging threats.

By integrating these concepts, the proposed system achieves a balance between
accuracy, adaptability, and practical usability, advancing the capabilities of
modern intrusion detection systems.